(2 points + 4 points)
\begin{teilaufgaben}
\item
The function $u(x,t)$ is defined on the upper half plane
$\Omega = (-\infty, \infty) \times [0,\infty)$.

The function $u(x,t)$ satisfies in $\Omega$ the equation
\[
u_{x}(x,t) + u_{t}(x,t) = 0 
\]
and $u(x,0) = x$ on the boundary of $\Omega$. 

Determine an approximate value for $u(0,0.2)$.

Use the scheme of Lax-Wendroff with $\Delta x = 0.12$ and $\Delta t = 0.10.$

\item
The function $u(x,t)$ is defined on the strip
$\Omega = [-1, 1] \times [0,\infty)$.

The function $u(x,t)$ satisfies in $\Omega$ the equation
\[
u_{xx}(x,t) = u_{tt}(x,t)
\]
and  $u(-1,t) = u(1,t) = 0$ on the boundary of $\Omega$. 

The function $u(x,t)$ also satisfies the initial conditions
\[
u(x,0) = 1-x^4 \qquad\text{and}\qquad u_t(x,0) = 0.
\]

Determine approximate values for $u(-1/2,1/2)$, $u(0,1/2)$ and $u(1/2,1/2)$.

Use finite differences with $\Delta x = 1/2$ and $\Delta t = 1/4.$
\end{teilaufgaben}

\begin{loesung}
\begin{teilaufgaben}
\item
Set
\[
U_{k,j} = \tilde u(0.12 \cdot k , 0.10 \cdot j).
\]
Apply the Lax-Wendroff scheme with $r = 5/6$ to obtain:
\begin{equation}
U_{k,j+1}
=
55/72 \cdot U_{k-1,j} + 22/72 \cdot U_{k,j} - 5/72 \cdot U_{k+1,j}.
\tag{{\bf 1P}}
\end{equation}
Therefore:
\[
\tilde u(-0.12,0.1) = -0.22,
\qquad
\tilde u(0,0.1) = -0.1,
\qquad
\tilde u(0.12,0.1) = 0.02
\]
and
\begin{equation}
\tilde u(0,0.2) = -0.2.
\tag{{\bf 1P}}
\end{equation}

\item
We have $f(x) = 1-x^4$ and $g(x) = 0$.
Hence for the first time step we have by Taylor
\[
u(x,1/4) \approx (1-x^4) + (-12 \cdot x^2) \cdot 1/32.
\]
This leads to 
\begin{equation}
\tilde u(-1/2, 1/4)
=
27/32, \qquad \tilde u(0, 1/4) = 1, \qquad \tilde u(0, 1/4) = 27/32. \qquad
\tag{{\bf 1P}}
\end{equation}    
For the second time step we use the leapfrog scheme with $r = 1/2$.
In this way we obtain the formula for $\tilde u(k \cdot 1/2, 1/2)$
(with $k = -1,0,1$):
\begin{equation}
1/4 \cdot \tilde u((k-1) \cdot 1/2, 1/4)
+ 6/4 \cdot \tilde u(k \cdot 1/2,1/4)
+ 1/4 \cdot  \tilde u((k+1) \cdot 1/2, 1/4)
- \tilde u(k \cdot 1/2,0)
\tag{{\bf 1P}}
\end{equation}
Finally we have
\begin{equation}
\tilde u(-1/2, 1/2)
=
\tilde u(1/2, 1/2) = 37/64, \qquad
\tilde u(0, 1/2) = 59/64.
\tag{{\bf 2P}}
\end{equation}
\end{teilaufgaben}
\end{loesung}

